\subsection{Fixed-set hardness and a dominated-law bridge to full AA}\label{ian:gamma-hard-sector-bridge}

Status: proved: non-adaptive hardness is characterized exactly by fixed-set posterior moments, and a dominated-law Global hypothesis implies full AA with explicit polynomial exponents (Theorem~\ref{ian:dominated-global-full-aa}). That theorem is a correct implication with no known natural instance: its hypothesis already fails for uniform product laws at \(w=1\), even for acceptance functions (addressing, Lemma~\ref{ian:addressing-reveal}), and every natural class examined in \path{ian/25_dominated_gamma_examples.tex} fails it.

Throughout, \(Q:\{-1,1\}^n\to[0,1]\), and \(h_S\), \(G_\nu(Q,S)\) are the fixed-set
posterior quantities of Definition~\ref{ian:adaptive-fixed-set-quantities} in
\path{ian/29_adaptive_transcript_hiding.tex}:
\[
 h_S=\E_\nu[Q\mid X_S]-\E_\nu Q,
 \qquad
 G_\nu(Q,S)=\E_\nu h_S^2.
\]
When the law and target must be displayed, write
\[
 R_\nu(Q,s)=\max_{|S|\le s}G_\nu(Q,S),
 \qquad
 M_\nu(Q,s)=\max_{|S|\le s}\E_\nu|h_S|,
\]
which are \(R_s\) and \(M_s\) of that definition. The adaptive test class
\(\mathcal F_t\), \(C_t(Q)=\sup_{F\in\mathcal F_t}|\Cov_\nu(Q,F)|\), and
\(\operatorname{Hard}\) are those of
Definition~\ref{ian:hard-sector-definition} in
\path{ian/13_hard_sector_definitions_forward.tex} and
Theorem~\ref{ian:adaptive-transcript-duality}; the influences
\(I_i^\nu\) and Gamma constants are those of
\path{ian/03_generalized_influence_bridge.tex}.

\begin{definition}[Non-adaptive hardness]
\label{ian:nonadaptive-hardness-definition}
For an integer \(t\ge0\), let
\[
 \mathcal F_t^{\mathrm{na}}
 =
 \bigcup_{\substack{S\subseteq[n]\\|S|\le t}}
 \{F(X_S):F:\{-1,1\}^S\to[-1,1]\}.
\]
Equivalently, these are trees that choose their queried coordinate set
before receiving any answers. Random mixtures of such tests do not
change the supremum of absolute covariance.

For \(0<\gamma<a\le1/4\), write
\(\operatorname{Hard}^{\mathrm{na}}(Q,\nu;t,a,\gamma)\) when
\[
 \Var_\nu Q\ge a,
 \qquad
 |\Cov_\nu(Q,F)|\le\gamma
 \quad(F\in\mathcal F_t^{\mathrm{na}}).
\]
\end{definition}

\begin{corollary}[Non-adaptive hardness in posterior moments]
\label{ian:fixed-set-hardness-duality}
For \(0<\gamma<a\le1/4\),
\[
 \operatorname{Hard}^{\mathrm{na}}(Q,\nu;t,a,\gamma)
 \quad\Longleftrightarrow\quad
 \Var_\nu Q\ge a,\quad M_\nu(Q,t)\le\gamma.
\]
In particular, assuming \(\Var_\nu Q\ge a\),
\[
 R_\nu(Q,t)\le\gamma^2
 \ \Longrightarrow\
 \operatorname{Hard}^{\mathrm{na}}(Q,\nu;t,a,\gamma)
 \ \Longrightarrow\
 R_\nu(Q,t)\le\gamma/2.
\]
\end{corollary}

\begin{proof}
Proposition~\ref{ian:adaptive-fixed-set-conversions} in
\path{ian/29_adaptive_transcript_hiding.tex} states that, for every fixed
\(S\), the supremum of \(|\Cov_\nu(Q,F(X_S))|\) over \(|F|\le1\) equals
\(\E_\nu|h_S|\), so that \(M_\nu(Q,t)\) is exactly the covariance
supremum over \(\mathcal F_t^{\mathrm{na}}\), and that
\(2R_\nu(Q,t)\le M_\nu(Q,t)\le\sqrt{R_\nu(Q,t)}\). This proves both
assertions. The lower comparison is attained by the affine test
\(2\E_\nu[Q\mid X_S]-1\in[-1,1]\), whose covariance with \(Q\) is
\(2G_\nu(Q,S)\); both constants are sharp by
Lemma~\ref{ian:adaptive-posterior-moments}.
\end{proof}

The conversions between fixed sets and adaptive trees are also in
Proposition~\ref{ian:adaptive-fixed-set-conversions}:
\(\operatorname{Hard}(Q,\nu;t,a,\gamma)\) implies \(R_\nu(Q,t)\le\gamma/2\);
conversely, with \(\Var_\nu Q\ge a\), either
\(R_\nu(Q,\min\{n,2^t-1\})\le\gamma^2\) or \(R_\nu(Q,t)\le2^{-t}\gamma^2\)
implies \(\operatorname{Hard}(Q,\nu;t,a,\gamma)\). The converses of
the first implication and of the condition
\(R_\nu(Q,\min\{n,2^t-1\})\le\gamma^2\) fail at unchanged parameters, by
Proposition~\ref{ian:adaptive-fixed-set-counterexamples}. The converse
of the condition \(R_\nu(Q,t)\le2^{-t}\gamma^2\) also fails: for \(n=3\),
the uniform law and \(Q=\frac12+\frac1{100}x_1+\frac15x_2x_3\), every
depth-one test has covariance at most \(1/100\), so
\(\operatorname{Hard}(Q,U;1,\frac{401}{10^4},\frac1{100})\) holds, while
\(R_U(Q,1)=10^{-4}>\frac12\cdot10^{-4}\).

\begin{remark}[What vanishing Local does and does not identify]
\label{ian:nonadaptive-hardness-versus-gamma-failure}
At fixed Gamma parameters, write \(\alpha\) for the Local threshold,
\(g\) for the Global threshold, and \(\ell\) for the allowed set size,
all applied to the \([0,1]\)-valued target \(Q\) (the affine
normalization \(Q=(1+f)/2\) divides influences and gains of \(f\) by
four). Failure of both \(\Gamma^{\mathrm{set}}\) alternatives for \(Q\)
under \(\nu\) means exactly
\[
 \max_i I_i^\nu(Q)<\alpha,
 \qquad
 R_\nu(Q,\ell)<g.
\]
Non-adaptive hardness instead controls the posterior \(L^1\) quantity
\(M_\nu(Q,\ell)\). The two are related through the transformations in
Corollary~\ref{ian:fixed-set-hardness-duality}; they are not identical
conditions at an unchanged threshold. For example, under the uniform
law on \(n=1\) coordinate, let \(0<b\le1/2\) and
\(Q=\frac12+bx_1\). Then \(h_{\{1\}}=bx_1\), so
\[
 R(Q,1)=b^2,\qquad M(Q,1)=b,
\]
and \(R(Q,1)<g<M(Q,1)\) for every \(g\in(b^2,b)\): the Global
alternative fails at threshold \(g\), while covariance tolerance \(g\)
fails.

More explicitly, Global failure at threshold \(g\) gives
\(M_\nu(Q,\ell)<\sqrt g\); with \(\Var_\nu Q\ge a\), this is legitimate
non-adaptive hardness with tolerance \(\sqrt g\) when \(\sqrt g<a\).
In the other direction the implication is one-sided and holds
literally: non-adaptive hardness with tolerance \(\gamma\) gives
\(R_\nu(Q,\ell)\le\gamma/2<\gamma\), so the Global alternative fails at
every threshold \(g>\gamma/2\), in particular at \(g=\gamma\).
In either direction, failure of the Local alternative is a separate
requirement. Exact vanishing \(I_i^\nu(Q)=0\) makes that requirement
automatic for every positive Local threshold.

Non-adaptive hardness does not imply small flip influences. Under
the uniform two-spin law,
\[
 Q=(1+X_1X_2)/2
\]
has variance \(1/4\), zero covariance with every depth-one test, and
\(I_1(Q)=I_2(Q)=1/4\).
\end{remark}

\begin{definition}[A dominated-law Global hypothesis]
\label{ian:dominated-global-hypothesis}
Let \(\mathcal A\) be a class of base laws. Fix constants
\[
 c,B,C>0,\qquad D,U,V,\lambda\ge0,\qquad 0\le\kappa<1.
\]
The dominated Global hypothesis means the following assertion,
uniformly over the base law, dimension, and all parameters.

For every \(n\ge1\), every \(\mu\in\mathcal A\) on \(\{-1,1\}^n\), every
integer \(d\ge1\), every \(0<w\le1\), every probability law
\(\nu\le\mu/w\), every globally bounded polynomial
\(Q:\{-1,1\}^n\to[0,1]\) of degree at most \(d\), and every
\(0<a\le1/4\), if \(\Var_\nu Q\ge a\), then there is \(S\subseteq[n]\)
such that
\[
 |S|\le
 \min\left\{n,\left\lceil B d^U a^{-V}w^{-\lambda}\right\rceil\right\},
 \qquad
 G_\nu(Q,S)\ge c\,d^{-D}a^Cw^\kappa.
\]
All constants precede \(\mu,n,d,Q,a,w,\nu\).

For a full-AA conclusion it suffices to require this assertion for
quantum acceptance functions, with \(d=\max\{1,2T\}\) for a
\(T\)-query acceptance function (the acceptance-function version).
Requiring it for all bounded polynomials is a stronger hypothesis.
The usual Gamma-shaped bounds correspond to
\[
 D=C,\qquad U=V=K,
\]
with a possible additional polynomial dependence on \(w^{-1}\).
The gain \(c\,a^Cw^\kappa\) without degree loss corresponds to \(D=0\).

Both versions are false for the class of uniform product laws, already
at \(w=1\): by Proposition~\ref{ian:dominated-global-product-failure} in
\path{ian/25_dominated_gamma_examples.tex}, the addressing functions of
Lemma~\ref{ian:addressing-reveal} are computed with \(r+1\) classical
queries but every fixed set has reveal gain at most \(|S|2^{-r}\).
\end{definition}

\begin{proposition}[Direct exclusion of an adaptive hard sector]
\label{ian:dominated-global-sector-exclusion}
Assume Definition~\ref{ian:dominated-global-hypothesis}. Let
\(\mu\in\mathcal A\) be a law on \(\{-1,1\}^n\), let \(d\ge1\), and let
\(Q:\{-1,1\}^n\to[0,1]\) be a globally bounded polynomial of degree at
most \(d\) (under the acceptance-function version, a \(T\)-query
acceptance function with \(d=\max\{1,2T\}\)). Let \(0<w\le1\) and
\(0<\gamma<a\le1/4\). No law \(\nu\le\mu/w\) satisfies
\(\operatorname{Hard}(Q,\nu;t,a,\gamma)\) if
\[
 t\ge\left\lceil B d^U a^{-V}w^{-\lambda}\right\rceil,
 \qquad
 c\,d^{-D}a^Cw^\kappa>\gamma/2.
\]
\end{proposition}

\begin{proof}
Hardness gives \(\Var_\nu Q\ge a\), so the hypothesis supplies \(S\)
with \(|S|\le t\) and \(G_\nu(Q,S)\ge c\,d^{-D}a^Cw^\kappa>\gamma/2\).
Proposition~\ref{ian:adaptive-fixed-set-conversions} gives
\(R_\nu(Q,t)\le\gamma/2\) under \(\operatorname{Hard}(Q,\nu;t,a,\gamma)\),
a contradiction.
\end{proof}

\begin{theorem}[Full AA from dominated Global progress]
\label{ian:dominated-global-full-aa}
Assume Definition~\ref{ian:dominated-global-hypothesis}. Put
\[
 s=\frac{D}{1-\kappa},
 \qquad
 k=\max\left\{0,\frac{2C+3\kappa-2}{1-\kappa}\right\},
 \qquad
 A=c\,4^{-C}16^{-\kappa},
\]
and choose
\[
 \zeta=\min\{1/4,A^{1/(1-\kappa)}\}.
\]
Define
\[
 P=U+s(3+\lambda),
 \qquad
 E=7+2V+3\lambda+k(3+\lambda),
\]
\[
 A_*=
 257\zeta^{-3}\left(1+B4^V16^\lambda \zeta^{-\lambda}\right).
\]
Let \(\mu\in\mathcal A\) be a law on \(\{-1,1\}^n\), let \(d\ge1\), and
let \(Q:\{-1,1\}^n\to[0,1]\) be a globally bounded polynomial of degree
at most \(d\) (under the acceptance-function version, a \(T\)-query
acceptance function with \(d=\max\{1,2T\}\)). For every
\(0<\varepsilon<1/8\), there is a deterministic adaptive tree with mean
absolute error at most \(\varepsilon\) under \(\mu\), of depth at most
\[
 \left\lceil A_*d^P\varepsilon^{-E}\right\rceil.
\]
Consequently \(\mathcal A\) has full AA uniformly over its base laws.
One admissible exponent in
\eqref{ian:hard-sector-aa-definition} is
\[
 B_0=\max\{1,\lceil P+E\rceil\}.
\]
For example, one may take
\[
 A_0=1+2^P16^{P+E}A_*.
\]
\end{theorem}

\begin{proof}
Fix \(d,Q,\varepsilon\), and set
\[
 \eta=\zeta d^{-s}\varepsilon^k,\qquad
 a=\frac{\varepsilon^2}{4},\qquad
 \gamma=\eta a,\qquad
 w=\frac{\eta\varepsilon^3}{16},
\]
\[
 t=\left\lceil B d^U a^{-V}w^{-\lambda}\right\rceil,
 \qquad
 L_*=\left\lceil256\varepsilon^{-7}\eta^{-3}\right\rceil t.
\]
Since \(0<\eta\le\zeta\le1/4\) and \(t\ge1\), these are admissible
parameters for Corollary~\ref{ian:hard-sector-explicit} in
\path{ian/13_hard_sector_definitions_forward.tex}, whose hypothesis
holds in particular when \(\ell_\mu(L_*,Q)>\varepsilon\).

Suppose \(\ell_\mu(L_*,Q)>\varepsilon\). That corollary supplies a set
sector \(\nu=\mu(\cdot\mid H)\le\mu/w\) satisfying
\[
 \operatorname{Hard}(Q,\nu;t,a,\gamma).
\]
The dominated Global gain is at least
\[
 g=c\,d^{-D}a^Cw^\kappa
   =A d^{-D}\eta^\kappa\varepsilon^{2C+3\kappa}.
\]
The ratio to its covariance tolerance is
\[
 \frac g\gamma
 =
 4A \zeta^{-(1-\kappa)}
 d^{-D+s(1-\kappa)}
 \varepsilon^{\,2C+3\kappa-2-k(1-\kappa)}
 \ge4.
\]
Here the degree exponent is zero, the epsilon exponent is
nonpositive, and \(\zeta^{1-\kappa}\le A\). In particular
\(g>\gamma/2\), and
Proposition~\ref{ian:dominated-global-sector-exclusion} gives a
contradiction. Thus \(\ell_\mu(L_*,Q)\le\varepsilon\).

For the depth calculation,
\[
 t\le
 \left(1+B4^V16^\lambda \zeta^{-\lambda}\right)
 d^{U+s\lambda}
 \varepsilon^{-(2V+3\lambda+k\lambda)},
\]
while
\[
 \left\lceil256\varepsilon^{-7}\eta^{-3}\right\rceil
 \le257\zeta^{-3}d^{3s}\varepsilon^{-(7+3k)}.
\]
Multiplication gives \(L_*\le A_*d^P\varepsilon^{-E}\), and monotonicity of
\(\ell_\mu(\cdot,Q)\) proves the depth bound.

For a \(T\)-query acceptance function, use
\(d=\max\{1,2T\}\), as in the proof of Theorem~\ref{ian:bridge}.
For other accuracies use
\(\varepsilon'=\min\{\varepsilon,1/16\}\). With
\(x=T+\varepsilon^{-1}\ge1\), one has \(d\le2x\) and
\(1/\varepsilon'\le16x\), so the depth is at most
\(\lceil A_*2^P16^Ex^{P+E}\rceil\le\lceil A_0x^{B_0}\rceil\).
\end{proof}

\begin{remark}[Bookkeeping, the common-exponent issue, and sector equivalence]
\label{ian:dominated-global-bookkeeping}
For the gain \(c\,a^Cw^\kappa\), one has \(s=0\), so
\(\eta=\zeta\varepsilon^k\) suffices. With a factor \(d^{-D}\) in the
gain, epsilon-only tuning does not uniformly absorb unbounded degree;
the factor \(d^{-D/(1-\kappa)}\) in \(\eta\) supplies what is missing.

For Gamma-shaped size and gain bounds,
\[
 P=K+\frac{C(3+\lambda)}{1-\kappa},
 \qquad
 E=7+2K+3\lambda
 +(3+\lambda)\max\left\{0,
       \frac{2C+3\kappa-2}{1-\kappa}\right\}.
\]
In particular, \(\kappa<1\) suffices even after the
\(\eta^{-3}\) factor in \(L_*\) is included. Since
Corollary~\ref{ian:hard-sector-explicit} already applies under
\(\ell_\mu(L_\dagger,Q)>\varepsilon\), with
\(L_\dagger=\lceil32\varepsilon^{-5}\eta^{-2}\rceil t\le
33\varepsilon^{-5}\eta^{-2}t\), the same proof also gives the smaller
exponents \(P=U+s(2+\lambda)\) and \(E=5+2V+3\lambda+k(2+\lambda)\), with
\(257\zeta^{-3}\) replaced by \(33\zeta^{-2}\) in \(A_*\).

The proof retains the actual extraction dispersion
\(a=\varepsilon^2/4\) and tolerance \(\gamma=\eta a\).
It never replaces \(a\) by a smaller power of epsilon.
Thus it does not use the invalid padding argument identified in
Remark~\ref{ian:hard-sector-common-exponent-warning}.

The canonical witnesses in
Theorem~\ref{ian:hard-sector-aa-equivalence} have \(\eta=1/4\).
Their gain need not contradict the dominated Global bound at that
fixed tolerance. Under full-AA failure, reapplying
Corollary~\ref{ian:hard-sector-explicit} with the displayed
degree-dependent and epsilon-dependent \(\eta\) gives the required
stronger hardness. Its inverse weight and extraction depth remain
fixed polynomials, chosen before any failure witnesses.

Equivalently, the theorem proves full AA in the form characterized
by Theorem~\ref{ian:hard-sector-aa-equivalence}. With the constants
\(A_0,B_0\) of the theorem, Lemma~\ref{ian:hard-sector-aa-exclusion} then
excludes, for every \(T\)-query acceptance function \(Q\), every smooth
sector of weight \(w\) with \(\operatorname{Hard}(Q,\nu;t,a,a/4)\) at a
depth
\[
 t\ge
 \left\lceil A_0\left(T+\frac8{3wa}\right)^{B_0}\right\rceil.
\]
The direct Global contradiction itself excludes the tuned sectors
above; it does not assert the absence of all hard sectors at every
depth.

At \(\kappa\ge1\), decreasing \(\eta\) no longer improves the ratio
\(g/\gamma\). The stated extraction calculation therefore supplies
no general extension to that range.
\end{remark}

\begin{remark}[Global progress versus a Local-or-Global assertion]
\label{ian:dominated-global-local-distinction}
The extracted sector is hard against adaptive tests, so it is
automatically hard against each permitted fixed-set test. Under
Definition~\ref{ian:dominated-global-hypothesis}, the contradiction
therefore uses Global progress only, and the proof needs no Local
alternative. This describes the proof, not the strength of the
hypothesis: dropping the Local alternative is exactly what makes the
hypothesis fail on natural classes. It is false for uniform product
laws at \(w=1\), also for acceptance functions, while those laws
satisfy \(\Gamma^{\mathrm{set}}\) through the Local alternative
(Proposition~\ref{ian:dominated-global-product-failure} and
Remark~\ref{ian:dominated-global-address-influence} in
\path{ian/25_dominated_gamma_examples.tex}). No natural class is known to
satisfy it.

A disjunction allowing also the Local alternative does not supply
the contradiction used here: adaptive hardness controls
correlations under the law, whereas flip derivatives can involve
off-support values. To see this, let
\(\operatorname{Hard}(Q,\nu;t,a,\gamma)\) hold with \(\gamma<a/2\). Add a
coordinate \(X_0\equiv1\) and define
\[
 Q'(x_0,x)=\frac14+\frac12Q(x)+\frac{x_0-1}{8}.
\]
This is globally in \([0,1]\). On the support, its variance and
covariances are respectively one quarter and one half of those of
\(Q\); querying \(X_0\) reveals no information, so
\(\operatorname{Hard}(Q',\delta_1\otimes\nu;t,a/4,\gamma/2)\) holds, which is
admissible because \(\gamma/2<a/4\). Meanwhile
\[
 I_0(Q')=1/64
\]
under every law. Hence, whenever the Local threshold at the new
parameters is at most \(1/64\), the Local alternative holds for this
hard target. Thus hardness alone cannot be used to discard a
flip-influence Local branch.

If a given target has \(I_i^\mu(Q)=0\) for every \(i\), this
vanishing survives every dominated law:
\[
 I_i^\nu(Q)\le I_i^\mu(Q)/w=0.
\]
For such targets a dominated Local-or-Global hypothesis forces
Global. For arbitrary targets, the argument above uses the Global
hypothesis; a Local-or-Global version would need an additional
argument handling Local progress.

The revealing-set budget must also fit inside the hardness depth:
\(t\ge\lceil B d^U a^{-V}w^{-\lambda}\rceil\).
An unspecified bound polynomial in the hardness depth \(t\) does
not itself guarantee this inequality.
\end{remark}

\begin{remark}[Relation to Theorem A]
\label{ian:dominated-global-versus-theorem-a}
Theorem~\ref{ian:dominated-global-full-aa} bypasses pinning closure
by applying Global progress directly to the extracted dominated
law. It proves full AA on the base class. It gives a hereditary
conclusion only if its hypothesis is also available uniformly for
the required pinned base laws. Its hypothesis has no known natural
instance: it fails for uniform products, triple blocks, Ising forests,
the character code and Curie--Weiss laws, and only an artificial
rare-sector family satisfies it
(\path{ian/25_dominated_gamma_examples.tex}). The theorem is therefore a
correct implication, not a source of new full-AA results.

For a base class \(\Lambda\) that is already pinning-closed, the
all-bounded-polynomial version of the hypothesis at \(w=1\) implies the
Global branch of \(\Gamma^{\mathrm{set}}(\Lambda)\), so
Theorem~\ref{ian:bridge} also applies. Indeed, for
\(f:\{-1,1\}^n\to[-1,1]\) of degree at most \(d\) with
\(\Var_\mu f\ge v\in(0,1]\), apply the hypothesis to \(Q=(1+f)/2\) with
\(a=v/4\); since \(G_\mu(f,S)=4G_\mu(Q,S)\), \(d\ge1\) and \(v\le1\), the
resulting set satisfies
\[
 |S|\le\min\left\{n,\left\lceil B4^V(d/v)^{\max\{U,V\}}\right\rceil\right\},
 \qquad
 G_\mu(f,S)\ge4^{1-C}c\,(v/d)^{\max\{C,D\}}.
\]
This is the Global alternative with \(B_G=B4^V\), \(K_G=\max\{U,V\}\),
\(c_G=4^{1-C}c\), \(C_G=\max\{C,D\}\), and arbitrary Local constants.
For an arbitrary base class \(\mathcal A\), the parameter-degraded
dominated hypothesis (gain allowed to decay like \(w^\kappa\)) and
\(\Gamma^{\mathrm{set}}(\operatorname{Pin}\mathcal A)\) are incomparable,
by Proposition~\ref{ian:dominated-global-pinning-incomparable} in
\path{ian/25_dominated_gamma_examples.tex}.
\end{remark}
